% Gerado por ferramentas/gerar.py; não edite à mão.
\section*{Exercícios}
\addcontentsline{toc}{section}{Exercícios}
Os pontos antes de cada número indicam a dificuldade: \nivel{1} fácil, \nivel{2} médio, \nivel{3} difícil e \nivel{4} desafio. Os exercícios marcados com \estrela{} têm solução completa, com animação do gráfico, no repositório do livro: \url{https://github.com/AleffAzevedo/calculo-1-pt/blob/main/exercicios/cap6/README.md}. As respostas estão no fim do capítulo.

\begin{itemize}[leftmargin=5.2em,labelsep=0.6em,itemsep=4pt]
\item[\exnum{1}{6.1}] Esboce a região entre $y=x^2$ e $y=2x-x^2$ e calcule a sua área.
\item[\exnum{1}{6.2}] Calcule a área entre $y=\sin x$ e $y=\cos x$ para $0\le x\le\tfrac\pi2$ (atenção ao cruzamento).
\item[\exnum{1}{6.3}] Calcule a área total da região limitada por $y=x^3$ e $y=x$.
\item[\exnum{1}{6.4}] Calcule a área entre $x=y^2$ e $x=4$ integrando em $y$.
\item[\exnum{1}{6.5}] Calcule a área entre $y=e^x$, $y=e^{-x}$ e a reta $x=1$.
\item[\exnum{1}{6.6}] Gire em torno do eixo $x$ a região sob $y=x^2$, $0\le x\le2$. Calcule o volume.
\item[\exnum{1}{6.7}] Gire em torno do eixo $x$ a região entre $y=\sqrt x$ e $y=x$. Calcule o volume (arruelas).
\item[\exnum{1}{6.8}] Gire em torno do eixo $y$ a região sob $y=\sqrt x$, $0\le x\le4$. Calcule o volume por cascas.
\item[\exnum{1}{6.9}] Deduza o volume de um cone circular reto de raio $r$ e altura $h$ girando uma reta adequada.
\item[\exnum{1}{6.10}] Calcule o comprimento de arco de $y=2x^{3/2}$ para $0\le x\le\tfrac13$.
\item[\exnum{1}{6.11}] São necessários $6$ J para esticar uma mola $0{,}3$ m além do comprimento natural. Quanto trabalho é preciso para esticá-la de $0{,}3$ m até $0{,}5$ m?
\item[\exnum{1}{6.12}] Uma barra ocupa $0\le x\le4$ (em metros) e tem densidade $\rho(x)=2+\sqrt x$ kg/m. Calcule a massa e diga se o centro de massa fica à esquerda ou à direita de $x=2$.
\item[\exnum{1}{6.13}] Calcule o centroide da região sob $y=2x$, $0\le x\le3$, e confira com o centroide do triângulo (média dos vértices).
\item[\exnum{1}{6.14}] Uma população de bactérias dobra a cada $3$ h. Quanto tempo leva para ficar $10$ vezes maior?
\item[\exnum{1}{6.15}] Um isótopo tem meia-vida de $8$ dias. De $200$ mg, quanto resta após $30$ dias?
\item[\exnum{1}{6.16}] Mostre que $\dfrac{d}{dx}(x\sinh x-\cosh x)=x\cosh x$ e use isso para calcular $\int_0^1x\cosh x\,dx$.
\item[\exnum{1}{6.17}] Calcule $\displaystyle\int\sinh x\cosh^2x\,dx$ e $\displaystyle\int_0^{\ln2}\cosh x\,dx$.
\item[\exnum{1}{6.18}] Usando $\ln2=\int_1^2\frac{dt}t$, estime $\ln2$ pela regra do ponto médio com $n=4$ e compare com $0{,}69315$.
\item[\exnum{1}{6.19}] Uma placa quadrada de $1$ m de lado está submersa na vertical, com a borda superior a $2$ m de profundidade. Calcule a força hidrostática numa face.
\item[\exnum{1}{6.20}] A base de um sólido é o disco $x^2+y^2\le4$; as seções perpendiculares ao eixo $x$ são quadrados. Calcule o volume.
\item[\exnum{2}{6.21}] \estrela\ \textbf{Área quando as curvas se cruzam.} Calcule a área total da região limitada por $y=x^3-3x$ e $y=x$.
\item[\exnum{2}{6.22}] Calcule a área entre $y=4-x^2$ e $y=x+2$.
\item[\exnum{2}{6.23}] Calcule a área da região limitada por $y=|x|$ e $y=2-x^2$.
\item[\exnum{2}{6.24}] Calcule a área da região limitada por $y=\ln x$, $y=0$ e $x=e$ integrando em $y$.
\item[\exnum{2}{6.25}] Calcule a área total entre $y=\sin x$ e $y=\cos x$ em $[0,2\pi]$.
\item[\exnum{2}{6.26}] Encontre $k$ tal que a reta $y=k$ divida a região entre $y=4-x^2$ e o eixo $x$ em duas partes de mesma área.
\item[\exnum{2}{6.27}] \estrela\ \textbf{Arruelas em torno de uma reta horizontal.} Gire a região entre $y=x$ e $y=x^2$ ($0\le x\le1$) em torno da reta $y=2$. Calcule o volume.
\item[\exnum{2}{6.28}] Gire em torno do eixo $x$ a região sob $y=\sec x$, $0\le x\le\tfrac\pi4$. Calcule o volume.
\item[\exnum{2}{6.29}] A base de um sólido é a região entre $y=x^2$ e $y=1$; as seções perpendiculares ao eixo $x$ são semicírculos. Calcule o volume.
\item[\exnum{2}{6.30}] \estrela\ \textbf{Cascas e arruelas: dois caminhos.} Gire a região entre $y=x$ e $y=x^2$ ($0\le x\le1$) em torno da reta $x=-1$. Calcule o volume pelo
método das cascas e confira com arruelas.
\item[\exnum{2}{6.31}] Gire em torno do eixo $y$ a região sob $y=e^{-x^2}$, $0\le x\le1$.
\item[\exnum{2}{6.32}] Para a região sob $y=\frac1x$, $1\le x\le3$, calcule o volume do giro em torno do eixo $x$ e em torno do eixo $y$.
\item[\exnum{2}{6.33}] \estrela\ \textbf{Um comprimento de arco que dá certo.} Calcule o comprimento de arco de $y=\dfrac{x^3}6+\dfrac1{2x}$ para $1\le x\le3$.
\item[\exnum{2}{6.34}] Calcule o comprimento de arco de $y=\ln(\cos x)$ para $0\le x\le\tfrac\pi4$.
\item[\exnum{2}{6.35}] Calcule a área da superfície obtida girando $y=x^3$, $0\le x\le1$, em torno do eixo $x$.
\item[\exnum{2}{6.36}] Uma corrente de $10$ m com $2$ kg/m pende do alto de um prédio. Quanto trabalho é preciso para içá-la toda até o topo?
\item[\exnum{2}{6.37}] Um disco de raio $2$ m tem densidade $\rho(r)=\frac1{1+r^2}$ kg/m² ($r$ = distância ao centro). Calcule a massa.
\item[\exnum{2}{6.38}] \estrela\ \textbf{Centroide e Pappus.} Encontre o centroide da região entre $y=x$ e $y=x^2$ e use o Teorema de Pappus para recalcular o
volume do giro em torno do eixo $x$ (Exemplo 6.3b).
\item[\exnum{2}{6.39}] \estrela\ \textbf{Crescimento de uma cultura.} Uma cultura de bactérias cresce a uma taxa proporcional ao seu tamanho e passa de $1000$ para $3000$
indivíduos em $2$ h. Encontre $k$, o tempo de duplicação, o tamanho após $5$ h e quando ela chega a $10\,000$.
\item[\exnum{2}{6.40}] R\$ $1000$ são aplicados a $8\%$ ao ano com juros compostos continuamente. Em quanto tempo o capital dobra? Quanto há após $10$ anos?
\item[\exnum{3}{6.41}] \estrela\ \textbf{Esvaziando um tanque cônico.} Um tanque cônico (vértice para baixo) tem $3$ m de raio no topo e $6$ m de altura e está cheio
de água. Calcule o trabalho para bombear toda a água até $1$ m acima do topo.
\item[\exnum{3}{6.42}] \estrela\ \textbf{Força numa comporta triangular.} Uma comporta vertical tem a forma de um triângulo com base de $4$ m na superfície da água e vértice
$3$ m abaixo. Calcule a força hidrostática sobre ela.
\item[\exnum{3}{6.43}] \estrela\ \textbf{O teorema de Arquimedes da faixa esférica.} Mostre que a área da faixa de uma esfera de raio $R$ entre dois planos paralelos a distância $h$
(ambos cortando a esfera) é $2\pi Rh$, qualquer que seja a posição dos planos. Que fração da superfície da
Terra fica acima da latitude $60^\circ$?
\item[\exnum{3}{6.44}] \estrela\ \textbf{O cabo suspenso.} Um cabo pende entre dois postes de mesma altura, separados por $40$ m, na forma $y=a\cosh\frac xa$
com $a=20$ m ($-20\le x\le20$). Mostre que o comprimento é $L=2a\sinh\frac{20}a$ e calcule $L$ e a flecha
(desnível entre os apoios e o ponto mais baixo).
\item[\exnum{3}{6.45}] \estrela\ \textbf{Volume do toro: Pappus × cascas.} Gire o disco $(x-3)^2+y^2\le1$ em torno do eixo $y$. Calcule o volume do toro por Pappus e
confira pelo método das cascas.
\item[\exnum{3}{6.46}] Calcule a área da região limitada por $y=x^2-2x$ e $y=4x-x^2$.
\item[\exnum{3}{6.47}] Calcule a área da região limitada por $y=\dfrac1{1+x^2}$ e $y=\dfrac{x^2}2$.
\item[\exnum{3}{6.48}] Gire em torno da reta $x=4$ a região limitada por $y=\sqrt x$, $y=0$ e $x=4$.
\item[\exnum{3}{6.49}] Gire em torno da reta $y=4$ a região limitada por $y=x^2$ e $y=4$.
\item[\exnum{3}{6.50}] Um furo cilíndrico é feito pelo centro de uma esfera de raio $R$, deixando um anel de altura $h$. Mostre que o volume que sobra é $\tfrac16\pi h^3$, independentemente de $R$.
\item[\exnum{3}{6.51}] Calcule o comprimento de arco de $y=\cosh x$ para $0\le x\le\ln2$.
\item[\exnum{3}{6.52}] Calcule o comprimento de arco de $y=\dfrac{x^2}4-\dfrac{\ln x}2$ para $1\le x\le e$.
\item[\exnum{3}{6.53}] Girando um segmento de reta, mostre que a área lateral de um cone de raio $r$ e geratriz $\ell$ é $\pi r\ell$.
\item[\exnum{3}{6.54}] Uma tigela hemisférica de raio $2$ m está cheia de água. Calcule o trabalho para bombear toda a água até a borda.
\item[\exnum{3}{6.55}] Uma janela circular de raio $0{,}5$ m está na parede vertical de um aquário, com o centro a $2$ m de profundidade. Calcule a força na janela por integral e verifique que $F=\rho g\,\bar h\,A$ ($\bar h$ = profundidade do centroide).
\item[\exnum{3}{6.56}] Use o Teorema de Pappus e o volume da esfera para mostrar que o centroide do semidisco de raio $r$ está a $\frac{4r}{3\pi}$ do diâmetro.
\item[\exnum{3}{6.57}] A força gravitacional sobre um satélite de $1000$ kg a distância $r$ do centro da Terra é $F=\frac{gR^2m}{r^2}$ ($R=6{,}4\times10^6$ m). Calcule o trabalho para elevá-lo da superfície até $1000$ km de altitude.
\item[\exnum{3}{6.58}] Pela lei de resfriamento de Newton, $T'=-k(T-20)$. Um café a $90$\textdegree C esfria para $60$\textdegree C em $10$ min. Quando chega a $40$\textdegree C? (Dica: $y=T-20$.)
\item[\exnum{3}{6.59}] Usando $\ln x=\int_1^x\frac{dt}t$, prove que $\ln x\le x-1$ para todo $x>0$ e conclua que $\tfrac12\le\ln2\le1$.
\item[\exnum{3}{6.60}] Calcule $\displaystyle\int\frac{dx}{\sqrt{x^2+4x+5}}$ completando o quadrado.
\item[\exnum{4}{6.61}] Um plano corta uma esfera de raio $R$. Mostre que uma calota de altura $h$ tem volume $\tfrac13\pi h^2(3R-h)$ e encontre $h$ para que o plano divida a esfera em volumes na razão $1:3$.
\item[\exnum{4}{6.62}] Mostre que para $y=\cosh x$ o comprimento de arco de $0$ a $b$ é igual à área sob a curva de $0$ a $b$. Prove que, entre as funções com $f(0)=1$ e $f'\ge0$, só $\cosh$ tem essa propriedade para todo $b$.
\item[\exnum{4}{6.63}] Calcule o volume comum a dois cilindros de raio $r$ cujos eixos se cruzam em ângulo reto (sólido de Steinmetz). (Dica: as seções paralelas aos dois eixos são quadrados.)
\item[\exnum{4}{6.64}] Encontre $f>0$ derivável com $f(0)=1$ tal que, para todo $x>0$: (a) a área sob $f$ de $0$ a $x$ seja $f(x)-1$; (b) essa área seja $f(x)^2-1$.
\item[\exnum{4}{6.65}] Prove que $\displaystyle\lim_{n\to\infty}\Big(\frac1{n+1}+\frac1{n+2}+\dots+\frac1{2n}\Big)=\ln2$.
\end{itemize}
\subsection*{Respostas}
{\small\begin{multicols}{2}\raggedright
\begin{itemize}[leftmargin=2.7em,labelsep=0.5em,itemsep=1pt]
\item[\textbf{6.1}] $\tfrac13$.
\item[\textbf{6.2}] $2\sqrt2-2\approx0{,}828$.
\item[\textbf{6.3}] $\tfrac12$.
\item[\textbf{6.4}] $\tfrac{32}3$.
\item[\textbf{6.5}] $e+e^{-1}-2\approx1{,}086$.
\item[\textbf{6.6}] $\tfrac{32\pi}5$.
\item[\textbf{6.7}] $\tfrac\pi6$.
\item[\textbf{6.8}] $\tfrac{128\pi}5$.
\item[\textbf{6.9}] Girando $y=\tfrac rhx$, $0\le x\le h$: $V=\tfrac13\pi r^2h$.
\item[\textbf{6.10}] $\tfrac{14}{27}$.
\item[\textbf{6.11}] $k=\tfrac{400}3$ N/m; $W=\tfrac{32}3\approx10{,}7$ J.
\item[\textbf{6.12}] $\tfrac{40}3$ kg; à direita (a barra é mais densa na ponta $x=4$).
\item[\textbf{6.13}] $(2,\,2)$.
\item[\textbf{6.14}] $3\tfrac{\ln10}{\ln2}\approx9{,}97$ h.
\item[\textbf{6.15}] $200\cdot2^{-30/8}\approx14{,}9$ mg.
\item[\textbf{6.16}] $1-e^{-1}\approx0{,}632$.
\item[\textbf{6.17}] $\tfrac13\cosh^3x+C$; $\tfrac34$.
\item[\textbf{6.18}] $M_4\approx0{,}69122$ (erro $\approx0{,}002$).
\item[\textbf{6.19}] $9800\int_2^3h\,dh=24\,500$ N.
\item[\textbf{6.20}] $\int_{-2}^24(4-x^2)\,dx=\tfrac{128}3$.
\item[\textbf{6.21}] $8$.
\item[\textbf{6.22}] $\tfrac92$.
\item[\textbf{6.23}] $\tfrac73$.
\item[\textbf{6.24}] $\int_0^1(e-e^y)\,dy=1$.
\item[\textbf{6.25}] $4\sqrt2$.
\item[\textbf{6.26}] $(4-k)^{3/2}=4$: $k=4-2^{4/3}\approx1{,}480$.
\item[\textbf{6.27}] $\tfrac{8\pi}{15}$.
\item[\textbf{6.28}] $\pi$.
\item[\textbf{6.29}] $\int_{-1}^1\tfrac\pi8(1-x^2)^2dx=\tfrac{2\pi}{15}$.
\item[\textbf{6.30}] $\tfrac\pi2$.
\item[\textbf{6.31}] $\pi(1-e^{-1})\approx1{,}986$.
\item[\textbf{6.32}] $\tfrac{2\pi}3$ e $4\pi$.
\item[\textbf{6.33}] $\tfrac{14}3$.
\item[\textbf{6.34}] $\int_0^{\pi/4}\sec x\,dx=\ln(1+\sqrt2)\approx0{,}881$.
\item[\textbf{6.35}] $\tfrac\pi{27}(10\sqrt{10}-1)\approx3{,}563$.
\item[\textbf{6.36}] $\int_0^{10}2\cdot9{,}8\,y\,dy=980$ J.
\item[\textbf{6.37}] $\int_0^2\tfrac{2\pi r}{1+r^2}dr=\pi\ln5\approx5{,}06$ kg.
\item[\textbf{6.38}] $(\bar x,\bar y)=\big(\tfrac12,\tfrac25\big)$; $V=\tfrac{2\pi}{15}$.
\item[\textbf{6.39}] $k=\tfrac{\ln3}2\approx0{,}549$/h; $T_2\approx1{,}26$ h; $\approx15\,588$; $\approx4{,}19$ h.
\item[\textbf{6.40}] $\tfrac{\ln2}{0{,}08}\approx8{,}66$ anos; $1000e^{0{,}8}\approx$ R\$ $2225{,}54$.
\item[\textbf{6.41}] $441\,000\pi\approx1{,}39\times10^6$ J.
\item[\textbf{6.42}] $58\,800$ N.
\item[\textbf{6.43}] $S=2\pi Rh$; cerca de $6{,}7\%$.
\item[\textbf{6.44}] $L=40\sinh1\approx47{,}0$ m; flecha $20(\cosh1-1)\approx10{,}9$ m.
\item[\textbf{6.45}] $6\pi^2\approx59{,}22$.
\item[\textbf{6.46}] $9$.
\item[\textbf{6.47}] $\tfrac\pi2-\tfrac13\approx1{,}237$.
\item[\textbf{6.48}] Cascas de raio $4-x$: $\tfrac{256\pi}{15}$.
\item[\textbf{6.49}] Discos de raio $4-x^2$: $\tfrac{512\pi}{15}$.
\item[\textbf{6.50}] Cascas: $V=2\pi\int_r^R2x\sqrt{R^2-x^2}\,dx=\tfrac43\pi(R^2-r^2)^{3/2}$ com $h=2\sqrt{R^2-r^2}$.
\item[\textbf{6.51}] $\sinh(\ln2)=\tfrac34$.
\item[\textbf{6.52}] $\tfrac{e^2+1}4\approx2{,}097$.
\item[\textbf{6.53}] $y=\tfrac rhx$: $S=2\pi\int_0^h\tfrac rhx\tfrac{\ell}{h}\,dx=\pi r\ell$.
\item[\textbf{6.54}] $9800\pi\int_0^2(4-u^2)u\,du=39\,200\pi\approx1{,}23\times10^5$ J.
\item[\textbf{6.55}] $4900\pi\approx15\,394$ N.
\item[\textbf{6.56}] $\tfrac43\pi r^3=2\pi\bar y\cdot\tfrac{\pi r^2}2$.
\item[\textbf{6.57}] $gmR\big(1-\tfrac R{R+h}\big)\approx8{,}48\times10^9$ J.
\item[\textbf{6.58}] $10\tfrac{\ln3{,}5}{\ln1{,}75}\approx22{,}4$ min.
\item[\textbf{6.59}] Compare $\tfrac1t$ com $1$ (e com $\tfrac12$ em $[1,2]$).
\item[\textbf{6.60}] $\sinh^{-1}(x+2)+C=\ln\big(x+2+\sqrt{x^2+4x+5}\big)+C$.
\item[\textbf{6.61}] $t=\tfrac hR$ satisfaz $t^3-3t^2+1=0$: $h\approx0{,}6527R$.
\item[\textbf{6.62}] Derivando: $\sqrt{1+f'^2}=f$, logo $f'=\sqrt{f^2-1}$, cuja solução com $f(0)=1$ é $\cosh x$.
\item[\textbf{6.63}] $\int_{-r}^r4(r^2-x^2)\,dx=\tfrac{16}3r^3$.
\item[\textbf{6.64}] (a) $f=e^x$; (b) $f=1+\tfrac x2$.
\item[\textbf{6.65}] É uma soma de Riemann de $\int_0^1\tfrac{dx}{1+x}=\ln2$.
\end{itemize}\end{multicols}}
